\section{Conjectures}\label{sec:conjectures}
This section contains no further theorem. The global gap rate in~\eqref{eq:rho} must be distinguished from the ratio between adjacent satellites at one rational scale.

\begin{conjecture}[The unresolved gap side]\label{conj:gap}
For $0<x=2\beta(\alpha)/L<1$,
\[
 \rho(\lambda,\alpha)=1-x.
\]
\end{conjecture}
The perturbative Mathieu formula at fixed $D$ suggests this value. It does not prove it. For $\kappa\gg1$, the $D$-th model gap is of order $\kappa^{1-D}$, much smaller than an error $o(\sqrt\kappa)$ when $D\ge2$. Resolving these gaps needs relative tunnelling estimates. An exact global rate also needs control over all sufficiently large labels. A local ladder supplies neither requirement.

\begin{conjecture}[The missing dynamical upper bound]\label{conj:transport}
If $1<x<2$, then
\[
 \alpha_u^{\rm unif}\le1-1/x.
\]
Together with the known lower bound this would give the exact worst exponent when $\beta_e=\beta$. A parity-free exact formula additionally requires a parity-free lower mechanism.
\end{conjecture}
A band-edge spectral comparison alone does not establish this conjecture. Position commutators, launch-state estimates, complementary-energy bounds and control between resonant scales are also needed.

\begin{conjecture}[Complementarity]\label{conj:complementarity}
The gap and worst-transport arithmetic laws are
\[
 \rho=(1-x)_+,\qquad
 \alpha_u^{\rm unif}=\alpha_{\rm worst}=(1-1/x)_+,
\]
with the appropriate endpoint interpretations. Equivalently, the proposed functions satisfy $\alpha_{\rm worst}(x)=\rho(1/x)$.
\end{conjecture}
The vanishing gap half for $x>1$ is now a theorem. The other equalities remain open in the ranges described above. Below one, zero fastest exponent is a weaker conjecture than dynamical localization. Literal all-phase dynamical localization is false: Jitomirskaya--Simon~\cite{JSimon} give dense sets of phases with no point spectrum. In the supercritical regime these have SC spectrum, precluding dynamical localization.

A stronger quantitative ladder beyond the second cap of~\eqref{eq:Dcap} is another open problem. A proof would have to retain both the inverse-polynomial function and the cosine potential inside a uniform quantization theorem. We do not import a general perturbation-stable Harper theorem under the name of the Mathieu ladder. Such a theorem would have additional parameter-uniformity obligations.

\label{end:conjectures}
